\section{Locally calibrated reconstruction from unlabelled counts}
\label{sec:v6-counts}
This section uses counts only, with neither channel labels nor impact
positions. There are three unoriented equal-gap ground channels. For an
actual-curvature conclusion we assume equal facing curvatures within
each channel. The asymmetry allowed in the general forward theorem is
not silently removed there: without the additional equality, the same
scalar reconstruction concerns \eqref{eq:v6-effective-curvature}.
All preparations are independent and have normalized full phase volume.
Collision counts and programmed physical times are recorded without
instrument noise. Statistical noise arises from the random preparation.

We consider either a sufficiently small compact subfamily of the
physical family in Theorem~\ref{thm:v6-three}, with $J=3$, or a compact
geometric family whose leading data lie in the four-amplitude neighborhood
of Theorem~\ref{thm:v5-pairwise}, with $J=4$. In the latter case the
normalizer is treated as an independent nuisance parameter in inversion;
the observed probabilities still come from a physical configuration.
For $1\le j\le J$, the ground-onset law reads
\begin{equation}\label{eq:v6-count-expansion}
 p_j(d):=\Pr\{N_{jg+d}\ge j+1\}
       =d^2\{C_j+dB_j(d)\},\qquad 0\le d\le d_*.
\end{equation}
The functions in braces are positive, bounded away from zero, and have
bounded derivatives at each fixed order, uniformly on the chosen family.
At these offsets the event is equivalently $N_{jg+d}=j+1$.

\begin{lemma}[Amplitude error from finitely many Bernoulli experiments]
\label{lem:v6-amplitudes}
Suppose $g$ is known. Fix $m\ge1$, and at each of the $Jm$ times
$jg+lh$, $1\le j\le J$, $1\le l\le m$, perform $n$ independent
preparations. Write $\widehat p_{j,l}$ for the observed success fraction
and put
\[
 \omega_l=(-1)^{l-1}\binom ml,\qquad
 \widehat C_j=\sum_{l=1}^m\omega_l
                         \frac{\widehat p_{j,l}}{(lh)^2}.
\]
For small $h$ and $0<\eta<1/2$, with probability at least $1-\eta$,
\begin{equation}\label{eq:v6-amplitude-concentration}
 \max_{j\le J}|\widehat C_j-C_j|
 \le C_m\left(h^m+
       \sqrt{\frac{\log(C_m/\eta)}{nh^2}}
                    +\frac{\log(C_m/\eta)}{nh^2}\right).
\end{equation}
\end{lemma}

\begin{proof}
The weights satisfy $\sum_l\omega_l=1$ and
$\sum_l\omega_l l^k=0$ for $1\le k<m$. Taylor's formula for
$p_j(d)/d^2$ in \eqref{eq:v6-count-expansion}, with its uniformly
bounded $m$th derivative, gives a population bias at most $C_mh^m$.
For a Bernoulli variable $X$ of mean $p$, exponential Markov bounds give
\[
 \Pr\{|\widehat p-p|>\sqrt{2px/n}+2x/(3n)\}\le2e^{-x}.
\]
For completeness, for $0<t<3$ the power series of the exponential and
$\E|X-p|^k\le\E(X-p)^2\le p$ for $k\ge2$ give
$\log\E e^{t(X-p)}\le pt^2/[2(1-t/3)]$; the same bound holds for
$-t$. Optimization for the sum of $n$ independent variables gives the
stated two-sided bound with these harmless enlarged constants.
Here $p_j(lh)\le C_m h^2$. Dividing by $(lh)^2$, summing the fixed
weights, and taking $x=\log(2Jm/\eta)$ proves the stochastic terms by
a union bound. Independence is used between preparations, not between
collisions of one trajectory.
\end{proof}

We now include fine timing instead of assuming it. A coarse bracket
$|g-g_0|\le h/4$ is supplied. Its acquisition is not charged below.
A pilot at order $j=1$ suffices for all the finitely many subsequent
orders. The safety margin in the second stage will be enlarged to
cover every possible pilot outcome.

\begin{theorem}[Count-only acquisition with charged fine calibration]
\label{thm:v6-count-only}
Fix $m\ge1$ and one of the preceding compact geometric families.
There are $h_m>0$ and $C_m<\infty$ such that, for
$0<h\le h_m$, $0<\eta<1/2$ and a supplied bracket
$|g-g_0|\le h/4$, an experiment using only unlabelled counts produces
estimates $\widehat g,\widehat A,\widehat\kappa$ with
\begin{equation}\label{eq:v6-count-errors}
 |\widehat g-g|\le C_mh^{m+1},\qquad
 |\widehat A-A|\le C_mh^m,\qquad
 \operatorname{dist}_{\rm match}(\widehat\kappa,\kappa)
                                         \le C_mh^{m/3}
\end{equation}
with probability at least $1-\eta$. Its total number $T$ of raw
preparations, including the pilot, satisfies
\begin{equation}\label{eq:v6-count-cost}
 \E T\le C_mh^{-(2m+2)}\log(C_m/\eta),\qquad
 T\le C_mh^{-(2m+2)}\log(C_m/\eta)
\end{equation}
with the second estimate holding jointly with
\eqref{eq:v6-count-errors} with probability at least $1-\eta$.
Consequently a sufficient raw-preparation order for matching curvature
accuracy $\varepsilon$ is
\begin{equation}\label{eq:v6-count-rate}
 C_m\varepsilon^{-(6+6/m)}\log(C_m/\eta).
\end{equation}
No channel labels, selected positions, or supplied fine-gap estimate
enter this experiment. The result is local and sufficient, not a
minimax theorem or a global procedure for locating the initial bracket.
\end{theorem}

\begin{proof}
Put $\tau=g-g_0$. At each of the pilot times $g_0+lh$,
$1\le l\le m+1$, count successes of $N_t\ge2$ until $r$ successes
have occurred, using fresh preparations at every window. Let $T_l$
be the waiting time and let $\widehat p_l=r/T_l$. Formula
\eqref{eq:v6-count-expansion} at $j=1$ has the form
\[
 p_1(d)=C_1d^2F(d),\qquad F(0)=1,
\]
where $F$ is positive and has the required bounded right derivatives.
The unknown $C_1>0$ replaces the unknown factor $K_j$ in
Lemma~\ref{lem:v5-calibration}; its proof uses no channel label.
Interpolate the square roots of the pilot fractions at $z=lh$ and
use precisely that lemma's root-or-zero rule. Thus
$\widehat g=g_0+\widehat\tau$ with
$|\widehat\tau|\le h/2$ on every outcome.
Choose
\[
 r=\left\lceil C_mh^{-2m}\log(C_m/\eta)\right\rceil.
\]
The negative-binomial bound \eqref{eq:v5-negative-binomial} gives
relative errors at most $c_mh^m$, simultaneously, outside an event of
probability $\eta/4$. The root estimate then gives
$|\widehat g-g|\le C_mh^{m+1}$.
Every pilot offset lies in $[3h/4,(m+5/4)h]$, so its success
probability is bounded below by $c_mh^2$. The expected pilot cost is
at most $C_mh^{-2}(m+1)r$. The same binomial waiting-time tail used
in Theorem~\ref{thm:v5-self-calibration} bounds the pilot cost by
$C_mh^{-(2m+2)}\log(C_m/\eta)$ except on a further event of probability
at most $\eta/4$.

For the second stage fix $L=J+1$ and perform $n$ fresh independent
preparations at each time $j\widehat g+Llh$, with $j\le J$, $l\le m$.
This stage observes only $\one_{\{N_t\ge j+1\}}$ and uses the estimator
\[
 \widehat C_j=\sum_{l=1}^m\omega_l
                   \frac{\widehat p_{j,l}}{(Llh)^2},\qquad
 n=\left\lceil C_mh^{-(2m+2)}\log(C_m/\eta)\right\rceil.
\]
The total second-stage cost $Jmn$ is deterministic. Even if the pilot
fails, $|\widehat g-g|\le3h/4$. Hence every actual second-stage offset
is between $(L-3J/4)h$ and $(Lm+3J/4)h$, and in particular exceeds
$h$. Decreasing $h_m$ places this entire range in the common collar.
Thus unsuccessful pilots cannot produce a nonpositive offset or an
undefined normalization.

Condition on the pilot. The fresh fractions have Bernoulli means
$p_j(Llh+j(\widehat g-g))$, each bounded by $C_mh^2$ uniformly over
all pilot outcomes. The concentration proof of
Lemma~\ref{lem:v6-amplitudes}, with nodes $Ll$, therefore still applies.
On the successful pilot event, differentiation of
\eqref{eq:v6-count-expansion} gives $|p_j'(d)|\le C_mh$ throughout the
relevant offset intervals. The mean-value theorem consequently gives
\[
 \left|\frac{p_j(Llh+j(\widehat g-g))-p_j(Llh)}{(Llh)^2}\right|
                 \le C_m\frac{|\widehat g-g|}{h}\le C_mh^m.
\]
This explicitly accounts for the uncancelled denominator/timing error.
Taylor extrapolation at the unshifted nodes $Llh$ has bias $C_mh^m$,
and the chosen $n$ makes both concentration terms at most $C_mh^m$.
A conditional union bound therefore gives
$\max_{j\le J}|\widehat C_j-C_j|\le C_mh^m$, with an additional
failure probability at most $\eta/2$.

In the three-amplitude physical case, fit the observations
$(\widehat g,\widehat C_1,\widehat C_2,\widehat C_3)$ over a sufficiently
small compact admissible family containing the true member and apply
Theorem~\ref{thm:v6-three}. Its true spacing discrepancy is
$O_m(h^{m+1})$ and its amplitude discrepancy is $O_m(h^m)$, so
\eqref{eq:v6-count-errors} follows, retaining the pilot as the stated
gap estimate. In the four-amplitude case, augment the coefficient map
in Theorem~\ref{thm:v5-pairwise} by the coordinate $g$. Its derivative
has the nonsingular four-by-four block already proved there, and the
identity in the $g$ coordinate. Thus its coefficient inverse is jointly
locally Lipschitz in $g$ and the four amplitudes. The same compact
minimum-discrepancy argument and the cubic root bound give the asserted
area and curvature estimates. On the failure set arbitrary compact
fallback estimates may be used; no success assertion depends on them.
Adding the pilot and deterministic second-stage costs proves
\eqref{eq:v6-count-cost}. Finally choose $h=c_m\varepsilon^{3/m}$,
with $c_m$ sufficiently small, to obtain \eqref{eq:v6-count-rate}.
\end{proof}

The exponent in \eqref{eq:v6-count-rate} has two different sources:
amplitude acquisition costs $h^{-(2m+2)}$, whereas coalescing unlabelled
curvatures have matching error $h^{m/3}$. The selected-position
experiment of Theorem~\ref{thm:v5-self-calibration} instead has a
Lipschitz curvature inverse and sufficient cost
$e^{j\gamma}\varepsilon^{-(2+2/m)}$. These are different observations
and different experiments. The sharp deterministic exponent $1/3$
does not by itself prove optimality of a statistical sampling rate.
All constants here may depend on fixed $m$, on $J\le4$, and on the
compact family; the statement makes no uniform extrapolation-order or
large-$j$ count-only cost claim.
